\chapter{Complexity — Deterministic}
%%% generated by the blueprint pipeline from TCSlib.Complexity.TuringMachine.Deterministic
%%%%%%%%%%%%%%%%%%%%%%%%%%%%%%%%%%%%%%%%%%%%%%%%%%%%%%%%%%%%%%%%%%%%%%%%%%%%%%%
\section{Overview}

This module defines deterministic multi-tape Turing machines with a read-only
input tape, $k$ work tapes, and a write-only output tape, together with their
step semantics: the one-step transition on configurations, the configuration
reached after $t$ steps, and the halting behaviour.  Space usage is measured
as the total number of work-tape cells visited by the heads, and on top of the
run semantics the module defines what it means for a machine to compute an
output, a function, or the indicator of a set within given time and space
bounds, including monotonicity of these notions in the bounds and uniqueness
of the halting time.

%%%%%%%%%%%%%%%%%%%%%%%%%%%%%%%%%%%%%%%%%%%%%%%%%%%%%%%%%%%%%%%%%%%%%%%%%%%%%%%
\section{Declarations}

\begin{definition}[Multi-tape Turing machine]
\lean{Turing.MultiTapeTM}
\label{Turing.MultiTapeTM}
\leanok
\uses{Turing.Action}
A deterministic Turing machine with a read-only input tape, $k$ work tapes,
and a write-only output tape, over an alphabet of symbols extended by a blank.
It consists of an initial state $q_0$ and a transition function that, given
the current state, the (possibly blank) symbol under the input head, and the
$k$ symbols under the work-tape heads, returns a one-step action: a movement
of the input head, a write-and-move instruction for each work tape, an
optional symbol to append to the output, and a successor state or the decision
to halt.  Neither the alphabet nor the state set is required to be finite.
\end{definition}

\begin{definition}[One computation step]
\lean{Turing.MultiTapeTM.step}
\label{Turing.MultiTapeTM.step}
\leanok
\uses{Turing.Action.apply, Turing.Cfg, Turing.Cfg.inputSymbol, Turing.Cfg.workTapeSymbols}
The successor of a configuration $c$ of a multi-tape Turing machine $M$: if
$c$ is halted (its state component is empty), the step returns $c$ unchanged;
otherwise it applies to $c$ the action that the transition function of $M$
prescribes for the current state, the symbol under the input head, and the
symbols under the $k$ work-tape heads.
\end{definition}

\begin{definition}[Symbol emitted by one step]
\lean{Turing.MultiTapeTM.outputSymbol}
\label{Turing.MultiTapeTM.outputSymbol}
\leanok
\uses{Turing.Cfg, Turing.Cfg.inputSymbol, Turing.Cfg.workTapeSymbols}
The symbol, if any, that a multi-tape Turing machine $M$ appends to its output
tape when executing one step from a configuration $c$: the optional output
component of the action prescribed by the transition function for the current
state and scanned symbols, and no symbol at all when $c$ is halted.
\end{definition}

\begin{definition}[Initial configuration]
\lean{Turing.MultiTapeTM.initCfg}
\label{Turing.MultiTapeTM.initCfg}
\leanok
\uses{Turing.Cfg, Turing.Cfg.init}
The starting configuration of a multi-tape Turing machine $M$ on an input
word: the machine is in its initial state $q_0$, the input head is at the
first input symbol, all work tapes are blank with heads at the origin, and the
output is empty.
\end{definition}

\begin{lemma}[A halted configuration is a fixed point of the step]
\lean{Turing.MultiTapeTM.step_of_halt}
\label{Turing.MultiTapeTM.step_of_halt}
\leanok
\uses{Turing.Cfg, Turing.MultiTapeTM.step}
\difficulty{1}
Let $M$ be a multi-tape Turing machine and $c$ a configuration whose state
component is empty (the machine has halted).  Then one computation step of $M$
leaves $c$ unchanged.
\end{lemma}

\begin{definition}[Configuration reached after $t$ steps]
\lean{Turing.MultiTapeTM.runFrom}
\label{Turing.MultiTapeTM.runFrom}
\leanok
\uses{Turing.Cfg, Turing.MultiTapeTM.step}
The configuration reached by iterating the one-step function of a multi-tape
Turing machine $t$ times starting from a configuration $c$.  Since halted
configurations are fixed points of the step function, a machine that halts
stays in its halting configuration.
\end{definition}

\begin{lemma}[Running zero steps]
\lean{Turing.MultiTapeTM.runFrom_zero}
\label{Turing.MultiTapeTM.runFrom_zero}
\leanok
\uses{Turing.Cfg, Turing.MultiTapeTM.runFrom}
\difficulty{1}
For every multi-tape Turing machine $M$ and configuration $c$, running $M$ for
$0$ steps from $c$ yields $c$ itself.
\end{lemma}

\begin{lemma}[Unfolding a run at the first step]
\lean{Turing.MultiTapeTM.runFrom_succ_eq_step}
\label{Turing.MultiTapeTM.runFrom_succ_eq_step}
\leanok
\uses{Turing.Cfg, Turing.MultiTapeTM.runFrom, Turing.MultiTapeTM.step}
\difficulty{1}
For every multi-tape Turing machine $M$, configuration $c$, and $t \in \mathbb{N}$,
running $M$ for $t + 1$ steps from $c$ equals running $M$ for $t$ steps from
the configuration obtained by one step of $M$ from $c$.
\end{lemma}

\begin{lemma}[Unfolding a run at the last step]
\lean{Turing.MultiTapeTM.runFrom_succ_eq_step'}
\label{Turing.MultiTapeTM.runFrom_succ_eq_step'}
\leanok
\uses{Turing.Cfg, Turing.MultiTapeTM.runFrom, Turing.MultiTapeTM.step}
\difficulty{1}
For every multi-tape Turing machine $M$, configuration $c$, and $t \in \mathbb{N}$,
running $M$ for $t + 1$ steps from $c$ equals applying one step of $M$ to the
configuration reached after running $M$ for $t$ steps from $c$.
\end{lemma}

\begin{lemma}[Splitting a run]
\lean{Turing.MultiTapeTM.runFrom_add}
\label{Turing.MultiTapeTM.runFrom_add}
\leanok
\uses{Turing.Cfg, Turing.MultiTapeTM.runFrom}
\difficulty{2}
For every multi-tape Turing machine $M$, configuration $c$, and
$a, b \in \mathbb{N}$, running $M$ for $a + b$ steps from $c$ equals running
$M$ for $b$ steps from the configuration reached after $a$ steps from $c$.
\end{lemma}

\begin{lemma}[A step-commuting map commutes with runs]
\lean{Turing.MultiTapeTM.runFrom_comm_of_step}
\label{Turing.MultiTapeTM.runFrom_comm_of_step}
\leanok
\uses{Turing.Cfg, Turing.MultiTapeTM, Turing.MultiTapeTM.runFrom, Turing.MultiTapeTM.step}
\difficulty{2}
Let $M$ and $M'$ be multi-tape Turing machines over the same alphabet, with
$k$ and $k'$ work tapes and running on inputs $x$ and $x'$ respectively, and
let $f$ map configurations of $M$ to configurations of $M'$ so that one step
of $M'$ after $f$ agrees with $f$ after one step of $M$ on every
configuration.  Then for every configuration $c$ of $M$ and every
$n \in \mathbb{N}$, running $M'$ for $n$ steps from $f(c)$ yields the image
under $f$ of the configuration $M$ reaches in $n$ steps from $c$.
\end{lemma}

\begin{lemma}[Runs from a halted configuration are constant]
\lean{Turing.MultiTapeTM.runFrom_of_halt}
\label{Turing.MultiTapeTM.runFrom_of_halt}
\leanok
\uses{Turing.Cfg, Turing.MultiTapeTM.runFrom, Turing.MultiTapeTM.step_of_halt}
\difficulty{1}
Let $M$ be a multi-tape Turing machine and $c$ a configuration whose state
component is empty (the machine has halted).  Then running $M$ for any number
of steps from $c$ yields $c$ itself.
\end{lemma}

\begin{lemma}[A halted configuration emits no output]
\lean{Turing.MultiTapeTM.outputSymbol_of_halt}
\label{Turing.MultiTapeTM.outputSymbol_of_halt}
\leanok
\uses{Turing.Cfg, Turing.MultiTapeTM.outputSymbol}
\difficulty{1}
Let $M$ be a multi-tape Turing machine and $c$ a configuration whose state
component is empty.  Then the step of $M$ from $c$ emits no output symbol.
\end{lemma}

\begin{lemma}[Work-tape heads move at most one cell per step]
\lean{Turing.MultiTapeTM.workTapePos_step_le}
\label{Turing.MultiTapeTM.workTapePos_step_le}
\leanok
\uses{Turing.Cfg, Turing.MultiTapeTM.step, Turing.workTapePos_apply_le}
\difficulty{2}
Let $M$ be a multi-tape Turing machine with $k$ work tapes, $c$ a
configuration, and $i$ one of the $k$ work-tape indices.  Then the position of
the $i$-th work-tape head after one step of $M$ from $c$ differs from its
position in $c$ by at most $1$ in absolute value.
\end{lemma}

\begin{definition}[Cells visited by a work-tape head]
\lean{Turing.MultiTapeTM.visitedByTapeHead}
\label{Turing.MultiTapeTM.visitedByTapeHead}
\leanok
\uses{Turing.Cfg, Turing.MultiTapeTM.runFrom}
The finite set of integer positions occupied by the head of work tape $i$
during the computation of a multi-tape Turing machine from a configuration
$c$, over steps $0$ through $t$ inclusive.
\end{definition}

\begin{definition}[Space used by one tape]
\lean{Turing.MultiTapeTM.spaceUsedByTape}
\label{Turing.MultiTapeTM.spaceUsedByTape}
\leanok
\uses{Turing.Cfg, Turing.MultiTapeTM.visitedByTapeHead}
The number of distinct work-tape cells touched by the head of work tape $i$ in
the computation of a multi-tape Turing machine from a configuration $c$ up to
step $t$, i.e.\ the cardinality of the set of positions visited by that head.
\end{definition}

\begin{definition}[Total space used]
\lean{Turing.MultiTapeTM.spaceUsed}
\label{Turing.MultiTapeTM.spaceUsed}
\leanok
\uses{Turing.Cfg, Turing.MultiTapeTM.spaceUsedByTape}
The total number of work-tape cells touched by the computation of a
multi-tape Turing machine with $k$ work tapes from a configuration $c$ up to
step $t$: the sum over the $k$ work tapes of the number of distinct cells
visited by each head.  This is the space measure used throughout; the
read-only input tape and the write-only output tape do not count.
\end{definition}

\begin{lemma}[A machine without work tapes uses no space]
\lean{Turing.MultiTapeTM.spaceUsed_zero_tapes_eq_zero}
\label{Turing.MultiTapeTM.spaceUsed_zero_tapes_eq_zero}
\leanok
\uses{Turing.Cfg, Turing.MultiTapeTM.spaceUsed}
\difficulty{2}
Let $M$ be a multi-tape Turing machine with $k$ work tapes where $k = 0$.
Then for every configuration $c$ and every step count $t$, the total space
used by $M$ from $c$ up to step $t$ is $0$.
\end{lemma}

\begin{lemma}[Per-tape space is at most the total space]
\lean{Turing.MultiTapeTM.spaceUsedByTape_le_spaceUsed}
\label{Turing.MultiTapeTM.spaceUsedByTape_le_spaceUsed}
\leanok
\uses{Turing.Cfg, Turing.MultiTapeTM.spaceUsed, Turing.MultiTapeTM.spaceUsedByTape}
\difficulty{1}
For every multi-tape Turing machine $M$, configuration $c$, step count $t$,
and work-tape index $i$, the number of cells touched by the head of tape $i$
up to step $t$ is at most the total space used by $M$ from $c$ up to step $t$.
\end{lemma}

\begin{lemma}[Space used equals the space of the computation trace]
\lean{Turing.MultiTapeTM.spaceUsed_eq_spaceUsedOfCfgs}
\label{Turing.MultiTapeTM.spaceUsed_eq_spaceUsedOfCfgs}
\leanok
\uses{Turing.Cfg, Turing.MultiTapeTM.runFrom, Turing.MultiTapeTM.spaceUsed, Turing.MultiTapeTM.spaceUsedByTape, Turing.MultiTapeTM.visitedByTapeHead, Turing.spaceUsedOfCfgs, Turing.visitedOfCfgs}
\difficulty{3}
Let $M$ be a multi-tape Turing machine, $c$ a configuration, and
$t \in \mathbb{N}$.  The total space used by $M$ from $c$ up to step $t$
equals the space touched by the list of configurations that $M$ traverses in
steps $0, 1, \dots, t$ from $c$, where the space touched by a list of
configurations is the sum over the work tapes of the number of distinct head
positions occurring in the list.
\end{lemma}

\begin{lemma}[One step appends the emitted symbol to the output]
\lean{Turing.MultiTapeTM.step_output}
\label{Turing.MultiTapeTM.step_output}
\leanok
\uses{Turing.Action.apply, Turing.Cfg, Turing.MultiTapeTM.outputSymbol, Turing.MultiTapeTM.step}
\difficulty{2}
For every multi-tape Turing machine $M$ and configuration $c$, the output
component after one step of $M$ from $c$ equals the output of $c$ followed by
the symbol emitted by that step, if any (and equals the output of $c$
unchanged when no symbol is emitted).
\end{lemma}

\begin{lemma}[The output is stable after halting]
\lean{Turing.MultiTapeTM.runFrom_output_eq_of_halt}
\label{Turing.MultiTapeTM.runFrom_output_eq_of_halt}
\leanok
\uses{Turing.Cfg, Turing.MultiTapeTM, Turing.MultiTapeTM.runFrom, Turing.MultiTapeTM.runFrom_add, Turing.MultiTapeTM.runFrom_of_halt}
\difficulty{3}
Let $M$ be a multi-tape Turing machine, $c$ a configuration, and
$\tau \le t$ natural numbers such that the configuration reached by $M$ after
$\tau$ steps from $c$ is halted.  Then the output after $t$ steps from $c$
equals the output after $\tau$ steps from $c$.
\end{lemma}

\begin{definition}[Computing an output in given time and space]
\lean{Turing.MultiTapeTM.ComputesInTimeAndSpace}
\label{Turing.MultiTapeTM.ComputesInTimeAndSpace}
\leanok
\uses{Turing.MultiTapeTM, Turing.MultiTapeTM.initCfg, Turing.MultiTapeTM.runFrom, Turing.MultiTapeTM.spaceUsed}
The proposition that a multi-tape Turing machine $M$, started in its initial
configuration on an input word, is halted after $t$ steps (hence halts within
$t$ steps, since halting is absorbing), that its output tape then holds
exactly the given output word, and that the computation up to step $t$ touches
exactly $s$ work-tape cells.  Neither the alphabet nor the state set is
required to be finite.
\end{definition}

\begin{definition}[Computing a function within input-indexed bounds]
\lean{Turing.MultiTapeTM.ComputesFunInTimeAndSpace}
\label{Turing.MultiTapeTM.ComputesFunInTimeAndSpace}
\leanok
\uses{Turing.MultiTapeTM, Turing.MultiTapeTM.ComputesInTimeAndSpace}
Given injective encodings of the domain and codomain into words over the
machine's alphabet, a multi-tape Turing machine $M$ computes a function $f$
within time bound $t$ and space bound $s$ (both functions of the input) if for
every input $a$ there are $t' \le t(a)$ and $s' \le s(a)$ such that $M$, on
the encoding of $a$, halts within $t'$ steps with output the encoding of
$f(a)$, touching exactly $s'$ work-tape cells.  The alphabet and state set
need not be finite.
\end{definition}

\begin{definition}[Computability within time and space bounds]
\lean{Turing.MultiTapeTM.ComputableInTimeAndSpace}
\label{Turing.MultiTapeTM.ComputableInTimeAndSpace}
\leanok
\uses{Turing.MultiTapeTM, Turing.MultiTapeTM.ComputesFunInTimeAndSpace}
A function $f$ is computable within input-indexed time bound $t$ and space
bound $s$, with respect to injective encodings of its domain and codomain into
binary words, if there exist a number $k$ of work tapes, a finite state set,
and a multi-tape Turing machine over the binary alphabet with those tapes and
states that computes $f$ between the given encodings within the bounds $t$
and $s$.
\end{definition}

\begin{definition}[Computability within bounds on the input length]
\lean{Turing.MultiTapeTM.ComputableInTimeAndSpaceOfLength}
\label{Turing.MultiTapeTM.ComputableInTimeAndSpaceOfLength}
\leanok
\uses{Turing.MultiTapeTM.ComputableInTimeAndSpace}
A function $f$ is computable within time bound $t$ and space bound $s$, where
$t, s : \mathbb{N} \to \mathbb{N}$ are functions of the encoded input length,
if there is a binary Turing machine with finitely many states that, for every
input $a$, computes the encoding of $f(a)$ from the encoding of $a$ in at most
$t(n)$ steps and touching at most $s(n)$ work-tape cells, where $n$ is the
length of the encoding of $a$.
\end{definition}

\begin{theorem}[Weakening the bounds of a computation]
\lean{Turing.MultiTapeTM.ComputesFunInTimeAndSpace.mono}
\label{Turing.MultiTapeTM.ComputesFunInTimeAndSpace.mono}
\leanok
\uses{Turing.MultiTapeTM, Turing.MultiTapeTM.ComputesFunInTimeAndSpace}
\difficulty{2}
Let $M$ be a multi-tape Turing machine that computes a function $f$ between
given injective encodings within input-indexed time bound $t$ and space bound
$s$, and let $t'$ and $s'$ be bounds with $t(a) \le t'(a)$ and
$s(a) \le s'(a)$ for every input $a$.  Then $M$ also computes $f$ between the
same encodings within the bounds $t'$ and $s'$.
\end{theorem}

\begin{theorem}[Computability is monotone in the bounds]
\lean{Turing.MultiTapeTM.ComputableInTimeAndSpace.mono}
\label{Turing.MultiTapeTM.ComputableInTimeAndSpace.mono}
\leanok
\uses{Turing.MultiTapeTM.ComputableInTimeAndSpace, Turing.MultiTapeTM.ComputesFunInTimeAndSpace.mono, Turing.MultiTapeTM.indicator}
\difficulty{2}
Let $f$ be a function computable within input-indexed time bound $t$ and
space bound $s$ with respect to given binary encodings, and let $t'$ and $s'$
be bounds with $t(a) \le t'(a)$ and $s(a) \le s'(a)$ for every input $a$.
Then $f$ is computable within the bounds $t'$ and $s'$ with respect to the
same encodings.
\end{theorem}

\begin{definition}[Boolean indicator of a set]
\lean{Turing.MultiTapeTM.indicator}
\label{Turing.MultiTapeTM.indicator}
\leanok
The Boolean indicator function of a set $L \subseteq \alpha$: the function
mapping $x$ to true when $x \in L$ and to false otherwise.
\end{definition}

\begin{definition}[Decidability within time and space bounds]
\lean{Turing.MultiTapeTM.DecidableInTimeAndSpace}
\label{Turing.MultiTapeTM.DecidableInTimeAndSpace}
\leanok
\uses{Turing.MultiTapeTM.ComputableInTimeAndSpace, Turing.MultiTapeTM.indicator}
A set $L \subseteq \alpha$ is decidable within input-indexed time bound $t$
and space bound $s$, with respect to an injective binary encoding of $\alpha$,
if the Boolean indicator function of $L$ is computable within these bounds,
the one-bit answer being encoded as the single-symbol output word.
\end{definition}

\begin{definition}[Halting at exactly step $t$]
\lean{Turing.MultiTapeTM.haltsAtStep}
\label{Turing.MultiTapeTM.haltsAtStep}
\leanok
\uses{Turing.MultiTapeTM, Turing.MultiTapeTM.initCfg, Turing.MultiTapeTM.runFrom}
The Boolean predicate stating that a multi-tape Turing machine $M$ on a given
input halts after exactly $t$ steps: started from its initial configuration,
$M$ is halted after $t$ steps but not after $t - 1$ steps (truncated
subtraction).  In particular the predicate is never true at $t = 0$: a machine
must perform at least one step to halt.
\end{definition}

\begin{lemma}[Uniqueness of the halting step]
\lean{Turing.MultiTapeTM.halting_step_unique}
\label{Turing.MultiTapeTM.halting_step_unique}
\leanok
\uses{Turing.MultiTapeTM, Turing.MultiTapeTM.haltsAtStep, Turing.MultiTapeTM.initCfg, Turing.MultiTapeTM.runFrom, Turing.MultiTapeTM.runFrom_add, Turing.MultiTapeTM.runFrom_of_halt}
\difficulty{4}
Let $M$ be a multi-tape Turing machine and fix an input word.  If $M$ halts on
that input after exactly $t_1$ steps and also after exactly $t_2$ steps, then
$t_1 = t_2$.
\end{lemma}

\begin{lemma}[A repeated non-halting configuration precludes halting]
\lean{Turing.MultiTapeTM.not_halts_of_repeat_nonhalt}
\label{Turing.MultiTapeTM.not_halts_of_repeat_nonhalt}
\leanok
\uses{Turing.Cfg, Turing.MultiTapeTM.runFrom, Turing.MultiTapeTM.runFrom_add}
\difficulty{5}
Let $M$ be a multi-tape Turing machine and let $c$ be a configuration that is
not halted.  If running $M$ for $t + 1$ steps from $c$ returns to $c$ itself,
then the computation of $M$ from $c$ never halts: for every number of steps
$t'$, the configuration reached after $t'$ steps from $c$ is not halted.
\end{lemma}
